% GENERATED FILE. Edit canonical lessons, then run npm run generate:books.
% canonical-source-hash: fnv1a64:25e20a1a4523de47
% canonical-lessons: JA-C291-nakunaru, JA-C291-soushiki, JA-C291-kanashimu, JA-C291-haka, JA-C291-inoru, JA-R291-review-286-291

\chapter{Saying Goodbye}
\label{ch:ja-saying-goodbye}
\hlchaptermodality{291}

\begin{chapteropening}
\textbf{By the end of this chapter:} \emph{I can say who died and how the family mourned.}

Complete the last lesson of chapter 291: I can say who died and how the family mourned.
\end{chapteropening}

\section[\texorpdfstring{\ja{なくなる} (nakunaru)}{なくなる (nakunaru)}]{\ja{なくなる} (nakunaru) — to pass away}

\label{lesson:JA-C291-nakunaru}



\begin{quote}
Before the new one: say the Japanese for to see (someone) off, then the Japanese for new.
\end{quote}

\subsection*{You'll want to know: \ja{なくなる}}
\textbf{\ja{なくなる}} — \emph{nakunaru} — \textquotedblleft{}to pass away\textquotedblright{}.

\begin{quote}
\textbf{\ja{きょねん、そぼが} \ja{なくなりました。}} — \emph{kyonen, sobo ga nakunarimashita} — \textquotedblleft{}My grandmother passed away last year.\textquotedblright{}
\end{quote}

Polite \textbf{\ja{なくなります}}, \emph{nakunarimasu}; a godan verb. Past (polite): \textbf{\ja{そふが} \ja{なくなりました。}}, \emph{sofu ga nakunarimashita}.

\begin{cousinweb}
\textbf{\ja{なくなる}} (\emph{nakunaru}) is a native Japanese word (wago), inherited rather than borrowed from Chinese; its kanji was attached later as a kun'yomi (JMdict).
\end{cousinweb}

\begin{grammarlens}[title={What you've built}]
One more word for telling what happened.
\end{grammarlens}

\subsection*{Your turn}
\begin{itemize}
\raggedright
  \item \emph{Say it:} \emph{nakunaru}
  \item \emph{Say it:} \emph{nakunaru}, once more
  \item \emph{Say it:} \emph{kyonen, sobo ga nakunarimashita} — My grandmother passed away last year.
  \item \emph{Recall:} say the Japanese for to see (someone) off, then the Japanese for new, then say \emph{nakunaru} again
\end{itemize}

\begin{tcolorbox}[breakable,colback=teal!4,colframe=teal!35!black,title={Before you move on}]
What does \ja{なくなる} mean? (\textquotedblleft{}To pass away\textquotedblright{}.) Say it once more.
\end{tcolorbox}
\section[\texorpdfstring{\ja{そうしき} (sōshiki)}{そうしき (sōshiki)}]{\ja{そうしき} (sōshiki) — a funeral}

\label{lesson:JA-C291-soushiki}



\begin{quote}
Before the new one: say the Japanese for new, then the Japanese for to pass away.
\end{quote}

\subsection*{You'll want to know: \ja{そうしき}}
\textbf{\ja{そうしき}} — \emph{sōshiki} — \textquotedblleft{}a funeral\textquotedblright{}.

\begin{quote}
\textbf{\ja{あした} \ja{そうしきが} \ja{あります。}} — \emph{ashita sōshiki ga arimasu} — \textquotedblleft{}The funeral is tomorrow.\textquotedblright{}
\end{quote}

\begin{cousinweb}
\textbf{\ja{そうしき}} (\emph{sōshiki}) is Sino-Japanese (kango): two Chinese characters read with their on'yomi, readings borrowed from Chinese. They are \emph{sō}, \textquotedblleft{}bury\textquotedblright{} (Middle Chinese \emph{tsangH}) and \emph{shiki}, \textquotedblleft{}ceremony\textquotedblright{} (\emph{syik}) (KANJIDIC2; Middle Chinese in Baxter's transcription).
\end{cousinweb}

\begin{grammarlens}[title={What you've built}]
One more word for telling what happened.
\end{grammarlens}

\subsection*{Your turn}
\begin{itemize}
\raggedright
  \item \emph{Say it:} \emph{sōshiki}
  \item \emph{Say it:} \emph{sōshiki}, once more
  \item \emph{Say it:} \emph{ashita sōshiki ga arimasu} — The funeral is tomorrow.
  \item \emph{Recall:} say the Japanese for new, then the Japanese for to pass away, then say \emph{sōshiki} again
\end{itemize}

\begin{tcolorbox}[breakable,colback=teal!4,colframe=teal!35!black,title={Before you move on}]
What does \ja{そうしき} mean? (\textquotedblleft{}A funeral\textquotedblright{}.) Say it once more.
\end{tcolorbox}
\section[\texorpdfstring{\ja{かなしむ} (kanashimu)}{かなしむ (kanashimu)}]{\ja{かなしむ} (kanashimu) — to grieve; to mourn}

\label{lesson:JA-C291-kanashimu}



\begin{quote}
Before the new one: say the Japanese for to pass away, then the Japanese for a funeral.
\end{quote}

\subsection*{You'll want to know: \ja{かなしむ}}
\textbf{\ja{かなしむ}} — \emph{kanashimu} — \textquotedblleft{}to grieve; to mourn\textquotedblright{}.

\begin{quote}
\textbf{\ja{かぞくは} \ja{みんな} \ja{かなしんでいます。}} — \emph{kazoku wa minna kanashinde imasu} — \textquotedblleft{}The whole family is in mourning.\textquotedblright{}
\end{quote}

Polite \textbf{\ja{かなしみます}}, \emph{kanashimimasu}; a godan verb. Past (polite): \textbf{\ja{かぞくが} \ja{かなしみました。}}, \emph{kazoku ga kanashimimashita}.

\begin{cousinweb}
\textbf{\ja{かなしむ}} (\emph{kanashimu}) is a native Japanese word (wago), inherited rather than borrowed from Chinese; its kanji was attached later as a kun'yomi (JMdict).
\end{cousinweb}

\begin{grammarlens}[title={What you've built}]
One more word for telling what happened.
\end{grammarlens}

\subsection*{Your turn}
\begin{itemize}
\raggedright
  \item \emph{Say it:} \emph{kanashimu}
  \item \emph{Say it:} \emph{kanashimu}, once more
  \item \emph{Say it:} \emph{kazoku wa minna kanashinde imasu} — The whole family is in mourning.
  \item \emph{Recall:} say the Japanese for to pass away, then the Japanese for a funeral, then say \emph{kanashimu} again
\end{itemize}

\begin{tcolorbox}[breakable,colback=teal!4,colframe=teal!35!black,title={Before you move on}]
What does \ja{かなしむ} mean? (\textquotedblleft{}To grieve; to mourn\textquotedblright{}.) Say it once more.
\end{tcolorbox}
\section[\texorpdfstring{\ja{はか} (haka)}{はか (haka)}]{\ja{はか} (haka) — a grave}

\label{lesson:JA-C291-haka}



\begin{quote}
Before the new one: say the Japanese for a funeral, then the Japanese for to grieve.
\end{quote}

\subsection*{You'll want to know: \ja{はか}}
\textbf{\ja{はか}} — \emph{haka} — \textquotedblleft{}a grave\textquotedblright{}.

\begin{quote}
\textbf{\ja{まいとし} \ja{そふの} \ja{はかに} \ja{いきます。}} — \emph{maitoshi sofu no haka ni ikimasu} — \textquotedblleft{}Every year we visit grandfather's grave.\textquotedblright{}
\end{quote}

\begin{cousinweb}
\textbf{\ja{はか}} (\emph{haka}) is a native Japanese word (wago), inherited rather than borrowed from Chinese; its kanji was attached later as a kun'yomi (JMdict).
\end{cousinweb}

\begin{grammarlens}[title={What you've built}]
One more word for telling what happened.
\end{grammarlens}

\subsection*{Your turn}
\begin{itemize}
\raggedright
  \item \emph{Say it:} \emph{haka}
  \item \emph{Say it:} \emph{haka}, once more
  \item \emph{Say it:} \emph{maitoshi sofu no haka ni ikimasu} — Every year we visit grandfather's grave.
  \item \emph{Recall:} say the Japanese for a funeral, then the Japanese for to grieve, then say \emph{haka} again
\end{itemize}

\begin{tcolorbox}[breakable,colback=teal!4,colframe=teal!35!black,title={Before you move on}]
What does \ja{はか} mean? (\textquotedblleft{}A grave\textquotedblright{}.) Say it once more.
\end{tcolorbox}
\section[\texorpdfstring{\ja{いのる} (inoru)}{いのる (inoru)}]{\ja{いのる} (inoru) — to pray}

\label{lesson:JA-C291-inoru}



\begin{quote}
Before the new one: say the Japanese for to grieve, then the Japanese for a grave.
\end{quote}

\subsection*{You'll want to know: \ja{いのる}}
\textbf{\ja{いのる}} — \emph{inoru} — \textquotedblleft{}to pray\textquotedblright{}.

\begin{quote}
\textbf{\ja{じんじゃで} \ja{かぞくの} \ja{けんこうを} \ja{いのりました。}} — \emph{jinja de kazoku no kenkō o inorimashita} — \textquotedblleft{}At the shrine I prayed for my family's health.\textquotedblright{}
\end{quote}

Polite \textbf{\ja{いのります}}, \emph{inorimasu}; a godan verb. Past (polite): \textbf{\ja{ははが} \ja{いのりました。}}, \emph{haha ga inorimashita}.

\begin{cousinweb}
\textbf{\ja{いのる}} (\emph{inoru}) is a native Japanese word (wago), inherited rather than borrowed from Chinese; its kanji was attached later as a kun'yomi (JMdict).
\end{cousinweb}

\begin{grammarlens}[title={What you've built}]
One more word for telling what happened.
\end{grammarlens}

\subsection*{Your turn}
\begin{itemize}
\raggedright
  \item \emph{Say it:} \emph{inoru}
  \item \emph{Say it:} \emph{inoru}, once more
  \item \emph{Say it:} \emph{jinja de kazoku no kenkō o inorimashita} — At the shrine I prayed for my family's health.
  \item \emph{Recall:} say the Japanese for to grieve, then the Japanese for a grave, then say \emph{inoru} again
\end{itemize}

\begin{tcolorbox}[breakable,colback=teal!4,colframe=teal!35!black,title={Before you move on}]
What does \ja{いのる} mean? (\textquotedblleft{}To pray\textquotedblright{}.) Say it once more.
\end{tcolorbox}
\section[(dialogue)]{Looking back: from A Baby Is Born to Saying Goodbye}

\label{lesson:JA-R291-review-286-291}



\begin{quote}
Say the Japanese for a baby and to be born.
\end{quote}

\subsection*{Your turn}
Say these aloud:

\begin{itemize}
\raggedright
  \item a baby, to be born, childbirth, congratulations, a grandchild
  \item toddling, to sprout, a grandmother, a relative, to grow up
  \item entering school, a ceremony, to attend, lively, a celebration
  \item a bride, a bridegroom, to vow, to invite formally, a married couple
  \item to move (house), newly married, one's home town, to see (someone) off, new
  \item to pass away, a funeral, to grieve, a grave, to pray
\end{itemize}

\begin{tcolorbox}[breakable,colback=teal!4,colframe=teal!35!black,title={Before you move on}]
A baby? (\textbf{\ja{あかちゃん}}, \emph{akachan}.) To be born? (\textbf{\ja{うまれる}}, \emph{umareru}.) Childbirth? (\textbf{\ja{しゅっさん}}, \emph{shussan}.) Congratulations? (\textbf{\ja{おめでとう}}, \emph{omedetō}.) A grandchild? (\textbf{\ja{まご}}, \emph{mago}.) Toddling? (\textbf{\ja{よちよち}}, \emph{yochiyochi}.) To sprout? (\textbf{\ja{はえる}}, \emph{haeru}.) A grandmother? (\textbf{\ja{おばあさん}}, \emph{obāsan}.) A relative? (\textbf{\ja{しんせき}}, \emph{shinseki}.) To grow up? (\textbf{\ja{せいちょうする}}, \emph{seichō suru}.) Entering school? (\textbf{\ja{にゅうがく}}, \emph{nyūgaku}.) A ceremony? (\textbf{\ja{しき}}, \emph{shiki}.) To attend? (\textbf{\ja{しゅっせきする}}, \emph{shusseki suru}.) Lively? (\textbf{\ja{にぎやか}}, \emph{nigiyaka}.) A celebration? (\textbf{\ja{おいわい}}, \emph{oiwai}.) A bride? (\textbf{\ja{はなよめ}}, \emph{hanayome}.) A bridegroom? (\textbf{\ja{はなむこ}}, \emph{hanamuko}.) To vow? (\textbf{\ja{ちかう}}, \emph{chikau}.) To invite formally? (\textbf{\ja{しょうたいする}}, \emph{shōtai suru}.) A married couple? (\textbf{\ja{ふうふ}}, \emph{fūfu}.) To move (house)? (\textbf{\ja{うつる}}, \emph{utsuru}.) Newly married? (\textbf{\ja{しんこん}}, \emph{shinkon}.) One's home town? (\textbf{\ja{ふるさと}}, \emph{furusato}.) To see (someone) off? (\textbf{\ja{みおくる}}, \emph{miokuru}.) New? (\textbf{\ja{あたらしい}}, \emph{atarashii}.) To pass away? (\textbf{\ja{なくなる}}, \emph{nakunaru}.) A funeral? (\textbf{\ja{そうしき}}, \emph{sōshiki}.) To grieve? (\textbf{\ja{かなしむ}}, \emph{kanashimu}.) A grave? (\textbf{\ja{はか}}, \emph{haka}.) To pray? (\textbf{\ja{いのる}}, \emph{inoru}.)
\end{tcolorbox}
